\section{잉여류와 라그랑주 정리}
\label{sec:cosets-lagrange}

부분군 $H\le G$가 주어졌을 때 $G$를 $H$와 같은 크기의 조각들로 나눌 수 있다. 이 조각이 잉여류이다.

\begin{definition}[왼쪽과 오른쪽 잉여류]
$H\le G$와 $a\in G$에 대해
\[
  aH=\{ah:h\in H\},\qquad
  Ha=\{ha:h\in H\}
\]
를 각각 $a$가 대표하는 \emph{왼쪽 잉여류}와 \emph{오른쪽 잉여류}라고 한다.
\end{definition}

잉여류는 새로운 부분군일 필요가 없다. $a\notin H$이면 보통 $aH$는 항등원을 포함하지 않는다. 잉여류의 역할은 부분군을 평행이동하여 $G$를 같은 크기의 블록으로 나누는 것이다.

\begin{example}
덧셈군 $\Z$에서 부분군 $3\Z$의 잉여류는
\[
  3\Z,\qquad 1+3\Z,\qquad 2+3\Z
\]
의 세 개이다. 이는 정수를 $3$으로 나눈 나머지에 따른 분할과 정확히 같다.
\end{example}

\begin{lemma}[잉여류 판정]
$H\le G$와 $a,b\in G$에 대해 다음은 동치이다.
\begin{enumerate}[label=(\roman*)]
  \item $aH=bH$.
  \item $aH\cap bH\ne\varnothing$.
  \item $b^{-1}a\in H$.
\end{enumerate}
\end{lemma}

\begin{proof}
(i)$\Rightarrow$(ii)는 $H$가 공집합이 아니므로 자명하다. (ii)를 가정하고 $x\in aH\cap bH$라 하자. 어떤 $h_1,h_2\in H$에 대해 $x=ah_1=bh_2$이므로
\[
  b^{-1}a=h_2h_1^{-1}\in H.
\]
따라서 (iii)가 성립한다. 마지막으로 $b^{-1}a\in H$이면 $a=b(b^{-1}a)\in bH$이고, 따라서 $aH\subseteq bH$이다. 같은 논리를 역원에 적용하면 $bH\subseteq aH$이므로 $aH=bH$이다.
\end{proof}

\begin{corollary}
서로 다른 두 왼쪽 잉여류는 서로소이다. 따라서 왼쪽 잉여류들은 $G$를 분할한다.
\end{corollary}

\begin{proof}
두 잉여류가 교차하면 보조정리에 의해 서로 같다. 또한 각 $g\in G$는 $g=ge\in gH$이므로 적어도 하나의 잉여류에 속한다.
\end{proof}

\begin{proposition}
모든 왼쪽 잉여류 $aH$는 $H$와 같은 원소 개수를 갖는다.
\end{proposition}

\begin{proof}
함수
\[
  L_a:H\to aH,\qquad h\mapsto ah
\]
는 전사이며, $ah_1=ah_2$이면 소거법칙에 의해 $h_1=h_2$이므로 단사이다. 따라서 전단사이다.
\end{proof}

\begin{definition}[지표]
$H$의 서로 다른 왼쪽 잉여류의 개수를 $H$의 $G$에서의 \emph{지표}라고 하고
\[
  [G:H]
\]
로 쓴다.
\end{definition}

왼쪽 잉여류와 오른쪽 잉여류의 개수는 같다. 실제로 $aH\mapsto Ha^{-1}$은 두 잉여류 집합 사이의 전단사이다.

\begin{theorem}[라그랑주 정리]
$G$가 유한군이고 $H\le G$이면
\[
  |G|=|H|[G:H].
\]
특히 $|H|$는 $|G|$를 나눈다.
\end{theorem}

\begin{proofidea}
잉여류들이 $G$를 서로 겹치지 않게 분할하고, 각 잉여류가 $H$와 같은 크기라는 두 사실을 곱한다.
\end{proofidea}

\begin{proof}
$G$의 서로 다른 왼쪽 잉여류를 $a_1H,\dots,a_rH$라 하자. 여기서 $r=[G:H]$이다. 이들은 $G$의 분할이고 각각 $|H|$개의 원소를 가지므로
\[
  |G|=\sum_{i=1}^r|a_iH|=r|H|=[G:H]|H|.
\]
\end{proof}

\subsection{원소의 위수에 대한 결과}

\begin{corollary}
유한군 $G$의 모든 원소 $g$에 대해
\[
  \ord(g)\mid |G|.
\]
따라서
\[
  g^{|G|}=e.
\]
\end{corollary}

\begin{proof}
$\langle g\rangle\le G$이므로 라그랑주 정리에 의해
\[
  \ord(g)=|\langle g\rangle|\mid |G|.
\]
$|G|=k\ord(g)$로 쓰면 $g^{|G|}=(g^{\ord(g)})^k=e$이다.
\end{proof}

\begin{corollary}[소수 위수 군]
$|G|=p$가 소수이면 $G$는 순환군이다. 더 나아가 $g\ne e$인 모든 원소가 $G$를 생성한다.
\end{corollary}

\begin{proof}
$g\ne e$이면 $\ord(g)>1$이고 $\ord(g)\mid p$이다. $p$가 소수이므로 $\ord(g)=p$이고 $\langle g\rangle=G$이다.
\end{proof}

\begin{corollary}[오일러 정리]
$\gcd(a,n)=1$이면
\[
  a^{\varphi(n)}\equiv1\pmod n,
\]
여기서 $\varphi(n)=|(\Z/n\Z)^\times|$는 오일러 파이 함수이다.
\end{corollary}

\begin{proof}
$\overline a\in(\Z/n\Z)^\times$에 라그랑주 정리를 적용하면
\[
  \overline a^{\,|(\Z/n\Z)^\times|}=\overline1
\]
이다. 이를 합동식으로 읽으면 원하는 결과를 얻는다.
\end{proof}

\begin{corollary}[페르마의 소정리]
$p$가 소수이고 $p\nmid a$이면
\[
  a^{p-1}\equiv1\pmod p.
\]
모든 정수 $a$에 대해 $a^p\equiv a\pmod p$이다.
\end{corollary}

\begin{remark}
라그랑주 정리의 역은 일반적으로 거짓이다. $d\mid|G|$라고 해서 항상 위수 $d$인 부분군이 존재하는 것은 아니다. 이후 실로 정리는 소수 거듭제곱 차수의 부분군에 대해서는 매우 강한 존재 정리를 제공하지만, 모든 약수에 대해 존재를 보장하지는 않는다.
\end{remark}

\subsection{잉여류 계산 예제}

\begin{example}
$S_3$에서 $H=\langle(12)\rangle=\{e,(12)\}$라 하자. 왼쪽 잉여류는
\[
\begin{aligned}
 H&=\{e,(12)\},\\
 (13)H&=\{(13),(13)(12)\}=\{(13),(123)\},\\
 (23)H&=\{(23),(23)(12)\}=\{(23),(132)\}.
\end{aligned}
\]
따라서 $[S_3:H]=3$이고 $|S_3|=2\cdot3=6$이다.
\end{example}

이 예에서 오른쪽 잉여류는 왼쪽 잉여류와 일치하지 않는다. 이 차이가 다음 절의 정규부분군 개념을 요구한다.

\begin{checkpoint}
\begin{enumerate}[label=\arabic*.]
  \item $\Z$에서 $4\Z$의 잉여류를 모두 쓰라.
  \item 위수 $24$인 군에서 가능한 원소의 위수는 어떤 수들로 제한되는가?
  \item 위수 $15$인 군의 원소 $g$가 $g^6=e$를 만족한다면 $g$의 가능한 위수는 무엇인가?
  \item 라그랑주 정리는 부분군의 존재를 보장하는 정리인가, 존재하는 부분군의 크기를 제한하는 정리인가?
\end{enumerate}
\end{checkpoint}
